\section{Teori Informasi dan Entropi}
Teori informasi mempelajari kuantisasi, penyimpanan, transmisi, dan pengolahan informasi. Salah satu metrik terpenting dalam kriptografi adalah \textbf{entropi}.

\subsection{Konsep Entropi Shannon}
Entropi mengukur tingkat ketidakpastian atau jumlah rata-rata informasi di dalam sebuah pesan. Semakin acak suatu data, semakin tinggi entropinya. Dalam kriptografi, cipherteks yang baik harus memiliki entropi yang tinggi agar sulit dianalisis.

Rumus entropi Shannon untuk variabel acak $X$:
\[ H(X) = - \sum_{i=1}^{n} p(x_i) \log_2 p(x_i) \]
di mana:
\begin{itemize}
    \item $x_i$ adalah simbol ke-$i$ di dalam pesan.
    \item $p(x_i)$ adalah peluang kemunculan simbol $x_i$.
\end{itemize}

\subsection{Rentang Nilai Entropi}
Nilai entropi berkisar antara:
\[ 0 \le H(X) \le \log_2(n) \]
\begin{itemize}
    \item \textbf{Entropi Minimum (0 bit)}: Terjadi jika tidak ada ketidakpastian (hanya satu simbol yang muncul dengan peluang 1).
    \item \textbf{Entropi Maksimum ($\log_2 n$ bit)}: Terjadi jika semua simbol muncul dengan peluang yang sama ($1/n$).
\end{itemize}

\textbf{Contoh Kasus:}
Untuk teks yang hanya terdiri dari 26 huruf alfabet (A--Z), entropi maksimumnya adalah $\log_2 26 \approx 4,7004$ bit/karakter. Jika sebuah cipherteks memiliki entropi yang mendekati nilai ini, maka ia sangat sulit dipecahkan menggunakan analisis frekuensi.





